\section{Resultados}\label{sec:resultados}

Este capítulo apresenta os resultados da campanha experimental
descrita no Capítulo~\ref{sec:metodologia}, organizados em seis
subseções. A primeira oferece a visão consolidada da comparação
entre os três controladores nos seis cenários. A segunda
discute, cenário a cenário, o comportamento observado. A
terceira analisa o trade-off central que emerge entre métricas
de cauda e métricas de regime estacionário. A quarta apresenta
a análise de sensibilidade do controlador AIMD CUBIC a
variantes de projeto da curva de recuperação. A quinta resume
a avaliação subjetiva. A sexta enuncia as limitações
metodológicas reconhecidas.

\subsection{Visão Geral}

A Tabela~\ref{tab:results-full} sintetiza, para cada par
(controlador, cenário), as métricas agregadas definidas no
Capítulo~\ref{sec:metodologia}. A Figura~\ref{fig:summary-bars}
visualiza as principais delas em forma comparativa.

\begin{table*}[!t]
  \caption{Resultados agregados da campanha experimental.
  Valores em negrito indicam o melhor desempenho entre os três
  controladores no respectivo cenário; as colunas de
  decrementos, incrementos e quadros IDR registram atividade do
  controlador e não admitem ordenação direta. Para as
  avaliações subjetivas (escala 1~=~melhor, 5~=~pior), o
  destaque marca o controlador subjetivamente preferido,
  conforme menor soma das três notas.}
  \label{tab:results-full}
  \centering
  \footnotesize
  \begin{tabular}{cl rrrr rr rr r c}
    \toprule
    \textbf{Cenário} & \textbf{Controlador} &
    \textbf{Pico} & \textbf{p99} & \textbf{p95} & \textbf{p50} &
    \textbf{Decr.} & \textbf{Incr.} &
    \textbf{Stalls} & \textbf{T. stall} & \textbf{IDR} &
    \textbf{lat/qual/estab} \\
    & & (ms) & (ms) & (ms) & (ms) & & & & (ms) & & \\
    \midrule
    \multirow{3}{*}{A} & sem ABR     & 1854 & 1801 & 1723 & \textbf{4} & 0 & 0 & 43 & 49\,337 & 10 & 5 / 1 / 5 \\
                      & AIMD de base & 1786 & 1713 & \textbf{55} & \textbf{4} & 7 & 25 & 4 & 4408 & 35 & \textbf{3 / 1 / 3} \\
                      & AIMD CUBIC   & \textbf{1025} & \textbf{913} & 183 & 16 & 9 & 82 & \textbf{0} & \textbf{0} & 97 & 3 / 2 / 3 \\
    \midrule
    \multirow{3}{*}{B} & sem ABR     & 1551 & 1481 & 1323 & \textbf{5} & 0 & 0 & 50 & 43\,550 & 27 & 4 / 1 / 4 \\
                      & AIMD de base & 575 & 544 & 194 & \textbf{4} & 7 & 22 & \textbf{0} & \textbf{0} & 30 & \textbf{1 / 1 / 2} \\
                      & AIMD CUBIC   & \textbf{532} & \textbf{486} & \textbf{123} & 18 & 7 & 67 & \textbf{0} & \textbf{0} & 78 & 3 / 2 / 2 \\
    \midrule
    \multirow{3}{*}{C} & sem ABR     & 1565 & 1541 & 1481 & \textbf{6} & 0 & 0 & 107 & 87\,452 & 104 & 4 / 1 / 4 \\
                      & AIMD de base & 733 & \textbf{504} & \textbf{116} & \textbf{5} & 7 & 28 & \textbf{0} & \textbf{0} & 36 & \textbf{1 / 1 / 1} \\
                      & AIMD CUBIC   & \textbf{639} & 566 & 298 & 13 & 12 & 115 & \textbf{0} & \textbf{0} & 133 & 2 / 1 / 3 \\
    \midrule
    \multirow{3}{*}{D} & sem ABR     & 1583 & 1531 & 1442 & 16 & 0 & 0 & 103 & 91\,501 & 97 & 4 / 1 / 3 \\
                      & AIMD de base & 1156 & \textbf{877} & \textbf{117} & \textbf{5} & 7 & 31 & \textbf{0} & \textbf{0} & 39 & \textbf{2 / 1 / 2} \\
                      & AIMD CUBIC   & \textbf{998} & 979 & 522 & 12 & 16 & 161 & \textbf{0} & \textbf{0} & 186 & 3 / 1 / 3 \\
    \midrule
    \multirow{3}{*}{E} & sem ABR     & 3788 & 3630 & 3435 & 22 & 0 & 0 & 59 & 135\,167 & 5 & 5 / 1 / 5 \\
                      & AIMD de base & 3808 & 3533 & 3354 & \textbf{5} & 16 & 25 & 5 & 11\,010 & 43 & 3 / 3 / 3 \\
                      & AIMD CUBIC   & \textbf{3311} & \textbf{3185} & \textbf{2980} & 23 & 16 & 128 & \textbf{2} & \textbf{4306} & 151 & 3 / 3 / 4 \\
    \midrule
    \multirow{3}{*}{F} & sem ABR     & 2536 & 2453 & 2308 & 22 & 0 & 0 & 100 & 148\,855 & 72 & 5 / 1 / 4 \\
                      & AIMD de base & \textbf{1602} & \textbf{1551} & \textbf{263} & \textbf{5} & 13 & 33 & \textbf{2} & \textbf{2852} & 50 & 3 / 2 / 3 \\
                      & AIMD CUBIC   & 1863 & 1832 & 477 & 12 & 17 & 128 & 4 & 5961 & 156 & \textbf{2 / 1 / 2} \\
    \bottomrule
  \end{tabular}
\end{table*}

\begin{figure*}[!t]
  \centering
  \includegraphics[width=\textwidth]{images/summary_bars.png}
  \caption{Comparação entre os três controladores nos seis
  cenários, agrupada por métrica. Da esquerda para a direita:
  pico de RTT, percentil 99, percentil 95, tempo total
  acumulado em \emph{stall}, e contagem de quadros IDR
  emitidos.}
  \label{fig:summary-bars}
\end{figure*}

Três observações orientam a discussão subsequente. Primeiro, o
controlador AIMD CUBIC apresenta o menor pico de RTT em cinco
dos seis cenários (A, B, C, D, E), com a maior redução
ocorrendo no cenário~A (1786 para 1025~ms, queda de 43\%).
Segundo, o mesmo controlador eleva o RTT mediano (p50) em
todos os cenários, passando de 4--5~ms no controlador AIMD de
base para 12--23~ms no AIMD CUBIC. Terceiro, a comparação com
a referência sem ABR é categórica: o tempo total acumulado em
\emph{stall} cai entre 90\% e 100\% nos dois controladores
adaptativos relativamente ao caso sem controle, com
\SI{40}{\second}--\SI{150}{\second} de \emph{stall} por sessão
de quatro minutos no controlador fixo. A ausência de ABR é
consistentemente a pior experiência, independentemente da
forma da degradação.

\subsection{Análise por Cenário}

A análise cenário a cenário esclarece em quais regimes cada
controlador opera melhor. A Figura~\ref{fig:timelines} reúne
as trajetórias de RTT ao longo das sessões, com o perfil de
\emph{shaping} sobreposto como referência visual.

\begin{figure*}[!t]
  \centering
  \includegraphics[width=\textwidth]{images/timelines.png}
  \caption{Trajetórias de RTT ao longo das sessões nos seis
  cenários, para cada controlador. As regiões sombreadas
  indicam as janelas de \emph{shaping} ativas. A escala de
  tempo é a duração da sessão (~\SI{240}{\second}); a escala de
  RTT é logarítmica para preservar a visibilidade simultânea
  do regime estacionário e dos picos.}
  \label{fig:timelines}
\end{figure*}

\textbf{Cenário A (queda abrupta).} A queda súbita de banda
livre para \SI{8}{\mega\bit\per\second} em $t=\SI{80}{\second}$
expõe o ganho mais nítido do controlador AIMD CUBIC. O pico de
RTT cai 43\% em relação ao AIMD de base (1025 \emph{vs.}
1786~ms), e o número de \emph{stalls} é reduzido a zero. A
combinação da curva cúbica de recuperação com o mecanismo de
\emph{fast convergence} absorve a restrição sem que o
codificador chegue a congelar. O controlador AIMD de base,
ainda que apresente p95 substancialmente menor (55 \emph{vs.}
183~ms), exibe um pico de cauda próximo do observado sem ABR.

\textbf{Cenário B (degraus sucessivos).} A redução gradual em
três degraus (16~/~12~/~8~Mbps) é o cenário em que o
controlador AIMD de base mais se beneficia de sua simplicidade.
A diferença no pico de RTT entre os dois controladores
adaptativos é pequena (575 \emph{vs.} 532~ms, queda de 7\% no
AIMD CUBIC). O regime estacionário, porém, é claramente
superior no AIMD de base, com p50 de 4~ms contra 18~ms do AIMD
CUBIC, e essa vantagem se reflete na avaliação subjetiva
amplamente favorável (1~/~1~/~2 \emph{vs.} 3~/~2~/~2). O p95
do AIMD CUBIC (123~ms) é, contudo, menor que o do AIMD de base
(194~ms), indicando que a contenção de picos isolados ocorre
mesmo nesse cenário em que o regime estacionário é a dimensão
dominante da percepção.

\textbf{Cenário C (dois patamares).} A transição lenta entre
14 e \SI{10}{\mega\bit\per\second} expõe um padrão similar ao
do cenário~B. O AIMD CUBIC reduz o pico em 13\%, mas o AIMD
de base alcança a única avaliação subjetiva perfeita da
campanha (1~/~1~/~1). O p95 do AIMD CUBIC (298~ms) é mais que
o dobro do AIMD de base (116~ms), e essa elevação do regime
estacionário é percebida pelo operador como degradação visual
ainda que os eventos de cauda sejam menos frequentes.

\textbf{Cenário D (janelas curtas alternadas).} O padrão de
quatro janelas de \SI{30}{\second} de constrição alternadas
com períodos livres é o que mais desafia o AIMD CUBIC. A cada
janela livre, o controlador inicia uma sondagem cúbica
agressiva; a cada nova constrição, dispara uma cascata de
cortes. O resultado é um p95 explodido (522 \emph{vs.} 117~ms
do AIMD de base) e uma sensação de instabilidade,
caracterizada pelas 161 atuações de incremento ao longo da
sessão. O pico de RTT é levemente menor no AIMD CUBIC (998
\emph{vs.} 1156~ms), mas a vantagem é eclipsada pela piora
do regime estacionário.

\textbf{Cenário E (constrição pesada longa).} A restrição
para \SI{4}{\mega\bit\per\second} sustentada por
\SI{160}{\second} é o cenário mais severo da campanha, e
aquele em que o argumento favorável ao AIMD CUBIC se
sustenta com maior força. O pico de RTT cai de 3808 para
3311~ms (13\%), e o tempo total em \emph{stall} é reduzido de
\SI{11.0}{\second} para \SI{4.3}{\second}, queda de 61\%. O
controlador sem ABR sofre \SI{135}{\second} de \emph{stall}
acumulado sobre os \SI{160}{\second} de constrição, equivalente
a 84\% do período: do ponto de vista do jogador, a sessão é
efetivamente tornada injogável.

\textbf{Cenário F (constrição moderada contínua).} A
restrição para \SI{6}{\mega\bit\per\second} ativa desde
$t=0$ é o único cenário em que o AIMD CUBIC regride
relativamente ao AIMD de base. O pico de RTT sobe 16\%
(1602 para 1863~ms) e o tempo em \emph{stall} é mais que
dobrado. O diagnóstico é estrutural: como a constrição existe
desde o início da sessão, o controlador nunca chega a operar
em uma fase não restrita; o teto de recuperação
\texttt{last\_known\_good} é inicializado com o valor
negociado de sessão (próximo de \SI{19}{\mega\bit\per\second})
e a curva inicial de recuperação fica apontada para esse teto
fictício, jamais alcançável. O mecanismo de \emph{fast
convergence} corrige progressivamente o teto à medida que
cascatas sucessivas ocorrem, mas a primeira janela de
recuperação é desperdiçada buscando uma capacidade que nunca
esteve disponível. Vale notar, contudo, que a avaliação
subjetiva no cenário~F favorece o AIMD CUBIC (2~/~1~/~2
\emph{vs.} 3~/~2~/~3 do AIMD de base), evidenciando que pico
de RTT e percepção subjetiva não se alinham de forma
mecânica.

\subsection{Trade-off entre Métricas de Cauda e de Regime Estacionário}

A leitura conjunta dos resultados revela um trade-off
sistemático entre o desempenho na cauda da distribuição de
latência (pico, percentil 99, eventos de \emph{stall}) e o
desempenho no regime estacionário (mediana e percentis
intermediários). A Figura~\ref{fig:tradeoff-scatter} apresenta
esse trade-off em uma única imagem.

\begin{figure}[!t]
  \centering
  \includegraphics[width=\columnwidth]{images/tradeoff_scatter.png}
  \caption{Trade-off entre pico de RTT (cauda) e mediana de
  RTT (regime estacionário). Cada ponto representa um par
  (controlador, cenário). Eixos em escala logarítmica. O canto
  inferior-esquerdo representa o ideal (ambas as métricas
  baixas); o canto superior-direito, o pior caso.}
  \label{fig:tradeoff-scatter}
\end{figure}

O agrupamento dos pontos por controlador é nítido. O AIMD de
base concentra-se na faixa de mediana de 4--5~ms, mas seus
picos cobrem todo o eixo vertical, alcançando valores acima
de \SI{3.5}{\second} no cenário~E. O AIMD CUBIC concentra-se
na faixa de mediana de 12--23~ms, mas seus picos são
consistentemente menores no mesmo cenário. Nenhum dos
controladores adaptativos domina o outro em ambos os eixos
simultaneamente: cada um define um regime de operação
distinto, e a fronteira de Pareto é composta por pontos de
ambos.

A Figura~\ref{fig:rtt-cdf} corrobora essa leitura ao
apresentar a distribuição cumulativa empírica do RTT para
cada par. As distribuições do AIMD de base sobem rapidamente
para próximo de 1 em valores baixos de RTT, refletindo o
regime estacionário privilegiado, mas exibem caudas longas
que se estendem além de \SI{1.5}{\second}. As distribuições
do AIMD CUBIC sobem mais gradualmente até o patamar de
\SI{100}{\milli\second}, refletindo o piso mais alto do
regime estacionário, mas concentram-se mais cedo nos valores
máximos, refletindo a contenção da cauda. As curvas dos dois
controladores adaptativos cruzam-se em todos os cenários,
geralmente em torno de \SI{100}{\milli\second}.

\begin{figure*}[!t]
  \centering
  \includegraphics[width=\textwidth]{images/rtt_cdf_panel.png}
  \caption{Distribuição cumulativa empírica (CDF) do RTT por
  cenário e por controlador. A linha vertical pontilhada
  marca o limiar de \SI{100}{\milli\second}. Eixos em escala
  logarítmica.}
  \label{fig:rtt-cdf}
\end{figure*}

O AIMD CUBIC reduz eventos de cauda ao custo de elevar o piso
de latência do regime estacionário, comportamento decorrente
da combinação do mecanismo de \emph{fast convergence} com a
curva cúbica de recuperação amostrada a 1~Hz. O AIMD de base,
por sua vez, otimiza o regime estacionário ao custo de menos
proteção contra eventos de cauda. Ambos os comportamentos
constituem escolhas de prioridade defensáveis: o primeiro
alinha-se à filosofia explicitada por Alós et al.~\cite{alos2019}
na seção~6, em que perdas de pacote e picos de latência são
priorizados sobre a utilização máxima do canal; o segundo
beneficia o ``caso comum'', em que a constrição é leve ou de
curta duração e o p50 percebido domina a experiência. A
escolha entre eles é, portanto, uma decisão de projeto sobre
qual frente da distribuição o controlador deve proteger,
não um juízo absoluto sobre qualidade.

\subsection{Análise de Sensibilidade ao Design da Curva}\label{sec:sensibilidade}

A curva de recuperação cúbica adotada na contribuição
(Capítulo~\ref{sec:contribuicao}) foi escolhida após avaliação
empírica de duas variantes alternativas, ambas mais fiéis à
formulação original do algoritmo CUBIC~\cite{cubic}. A primeira
variante, identificada como~5b, ancora o tempo decorrido $t$ da
curva no instante de retorno do RTT à condição de via livre
(\texttt{rtt\_clear\_since}) em vez do instante da última redução
(\texttt{last\_decrease\_at}). A motivação era teórica: o $t-K$
da curva pretende medir o tempo desde que o canal está
genuinamente livre, e \texttt{rtt\_clear\_since} captura isso
melhor que \texttt{last\_decrease\_at}, contaminado pelo
tempo de drenagem do \emph{buffer}. A segunda variante,
identificada como~5c, mantém a âncora em
\texttt{last\_decrease\_at} mas deixa de reinicializar o
\emph{streak} de via livre após cada incremento bem-sucedido,
permitindo que $t$ acumule monotonicamente entre cortes. Essa
é a interpretação literal do parágrafo §3.1 do trabalho
original do CUBIC.

A Tabela~\ref{tab:sensitivity} apresenta o pico de RTT obtido
em sessões repetidas dos cenários~A e~E com cada variante,
mantidos os demais parâmetros.

\begin{table}[!t]
  \caption{Pico de RTT (em ms) sob três variantes de design da
  curva de recuperação. A forma final é a adotada na
  contribuição. Variantes regrediram em todos os casos
  avaliados.}
  \label{tab:sensitivity}
  \centering
  \begin{tabular}{lrrr}
    \toprule
    \textbf{Variante} & \textbf{A K=5\,s} & \textbf{A K=15\,s} & \textbf{E K=15\,s} \\
    \midrule
    Forma final (Idea 5)  & 796   & 1221  & \textbf{2305} \\
    Variante 5b           & 1105 \scriptsize($+39\%$) & 1313 \scriptsize($+8\%$)  & 3192 \scriptsize($+39\%$) \\
    Variante 5c           & 1116 \scriptsize($+40\%$) & 1264 \scriptsize($+4\%$)  & 3013 \scriptsize($+31\%$) \\
    \bottomrule
  \end{tabular}
\end{table}

A regressão é sistemática nos cenários constrangidos,
particularmente no cenário~E, em que ambas as variantes
elevam o pico em mais de 30\% relativamente à forma final. O
diagnóstico é estrutural. A curva cúbica do CUBIC original foi
formulada para amostragem por~ACK, regime em que cada
atualização da janela aplica um delta pequeno acumulado ao
longo de muitas iterações por segundo. No controlador aqui
proposto, a curva é amostrada a uma cadência de 1~Hz, em que
cada atualização aplica o delta integral entre dois instantes
de decisão. As variantes mais fiéis ao paper original mostram
\emph{overshoot} sistemático: o controlador cresce mais do que
o canal pode absorver, dispara congestionamento novamente, e
permanece em ciclo. A variante simplificada, ao reinicializar
o \emph{streak} de via livre após cada crescimento, impõe um
intervalo mínimo de espera entre crescimentos sucessivos, o
que funciona como uma forma implícita de filtragem temporal
ausente na formulação original.

Esse achado tem alcance maior que a presente contribuição.
Algoritmos de controle de congestionamento projetados para
amostragem densa não são, em geral, transponíveis a regimes de
cadência reduzida sem ajuste estrutural. No caso do CUBIC, a
adaptação que se mostrou empiricamente superior é também a
mais simples, contradizendo a intuição de que maior fidelidade
ao paper original implicaria melhor desempenho.

\subsection{Avaliação Subjetiva}

A Figura~\ref{fig:ratings} apresenta as três avaliações
subjetivas registradas ao final de cada sessão.

\begin{figure*}[!t]
  \centering
  \includegraphics[width=\textwidth]{images/ratings.png}
  \caption{Avaliações subjetivas (escala 1~=~melhor,
  5~=~pior) por cenário e por controlador, organizadas em
  três eixos: latência percebida, qualidade visual e
  estabilidade.}
  \label{fig:ratings}
\end{figure*}

A avaliação subjetiva contradiz parcialmente a ordenação
sugerida pelas métricas de cauda. O controlador AIMD de base
recebe avaliações estritamente melhores nos cenários~B, C e~D,
apesar de seu pico de RTT ser igual ou pior que o do AIMD
CUBIC nesses mesmos cenários. O AIMD CUBIC, por sua vez,
recebe avaliação subjetiva claramente superior no cenário~F,
em que perde nos picos de RTT. Nos cenários~A e~E as
avaliações são próximas. A leitura natural é que o regime
estacionário (refletido em p50 e p95) tem peso perceptual
maior que eventos de cauda isolados, exceto quando esses
eventos ocorrem em magnitude e frequência elevadas, como no
cenário~E.

Esse padrão alinha-se às observações de QoE em jogos em
nuvem discutidas no Capítulo~\ref{sec:referencial}: a
percepção do usuário é mais sensível à constância do que ao
pico absoluto, contanto que o pico permaneça abaixo de um
limiar de tolerância. Quando esse limiar é ultrapassado, como
no cenário~E, a avaliação subjetiva se aproxima à medida que
o desempenho objetivo se distingue.

\subsection{Limitações dos Resultados}

A campanha experimental opera sob restrições metodológicas
que devem ser explicitadas para correta interpretação dos
resultados. Em primeiro lugar, foi realizada uma única
execução por par (controlador, cenário) para o controlador
AIMD CUBIC, sem possibilidade de quantificação estatística
da variância entre execuções. As diferenças observadas entre
controladores podem, em certos pares, ter magnitude
comparável à variância natural entre repetições. Essa
limitação foi assumida deliberadamente, dado o custo de
operação manual de cada sessão e a necessidade de manter o
operador subjetivo consistente em todas elas.

Em segundo lugar, as colunas referentes aos controladores sem
ABR e AIMD de base foram reaproveitadas de campanhas
anteriores, conduzidas no mesmo \emph{testbed} e com o mesmo
operador, mas separadas no tempo do conjunto de execuções do
AIMD CUBIC. Embora a configuração do servidor tenha
permanecido inalterada, fatores ambientais como tráfego
concorrente na rede doméstica e estado do sistema operacional
podem ter sofrido deriva no intervalo. A baixa variância
observada nas campanhas originais sustenta o reaproveitamento
como hipótese de trabalho, mas não como certeza estatística.

Em terceiro lugar, o desempenho efetivo do codificador
depende do conteúdo de vídeo exibido em cada sessão. Em
duas execuções iniciais dos cenários~A e~B, o jogo passava
por uma sequência menos exigente do ponto de vista visual,
em que a taxa de bits demandada pelo codificador permanecia
abaixo da capacidade imposta pelo \emph{shaping}: o
controlador não foi acionado por falta de pressão sobre o
canal. Essas execuções foram descartadas em favor de novas
sessões com atividade de \emph{gameplay} equivalente entre
controladores. A constatação reforça que a comparação entre
controladores requer não apenas um cenário de rede comum,
mas também uma carga de conteúdo comum.

Em quarto lugar, as avaliações subjetivas foram conduzidas
por um único operador, o autor do trabalho. Essa configuração
oferece consistência entre sessões, mas não captura
variabilidade interpessoal nem fornece a base estatística
necessária para sustentar conclusões sobre QoE em sentido
amplo. Estudos de MOS \emph{crowdsourced} como o conduzido por
Hong et al.~\cite{hong2015} representam o protocolo apropriado
para essa finalidade, e ficam como direção natural de trabalho
futuro.

Em quinto lugar, a contagem de \emph{stalls} reportada na
configuração sem ABR reflete pausas no decodificador, não
desconexões da sessão. No cenário~E, em que o tempo
acumulado em \emph{stall} alcança \SI{135}{\second} sobre
\SI{160}{\second} de constrição, é provável que a sessão
tenha efetivamente perdido sincronia com o servidor por
parte do período; o instrumento de coleta atual não distingue
esses dois regimes, e a métrica deve ser lida como
``intervalos sem entrega de quadros'' e não ``conexão ativa
sem entrega de quadros''.
